% Situaciones de aprendizaje de cada unidad.
% Se incluye en el anexo de Libro_4_ESO_ES.tex. Ordenado por unidades.

%========================================================================================
\section[Unidad 1. El pensamiento científico]{Unidad 1. La evolución del pensamiento científico}
%========================================================================================
% (sacado del final de UD1_Trabajo_cientifico.tex)

Todo lo que sucede en la naturaleza se puede organizar y explicar con unas pocas leyes: fenómenos aparentemente distintos, como la caída de una manzana y la órbita de la Luna, se explican con la ley de gravitación universal. El conocimiento científico ha evolucionado a lo largo de la historia y, a menudo, ha estado influido por el contexto político, económico o religioso. La ciencia moderna es fruto del trabajo de personas como Galileo, Bacon, Descartes o Newton, que usaron el método científico: un procedimiento experimental, riguroso y cuantitativo.

\begin{marker}
    \textbf{Reto:} investigar la evolución del conocimiento científico sobre un tema, relacionando los descubrimientos con sus protagonistas y con su contexto histórico, y elaborar una \textbf{línea del tiempo} y una presentación.
\end{marker}

Temas posibles: el origen del universo y los métodos para observarlo; el movimiento de estrellas, planetas y satélites; la naturaleza de la luz; la estructura atómica; la radiactividad; la química del carbono; la meteorología; el Sistema Internacional de unidades; las máquinas térmicas; el calentamiento global y el cambio climático.

\begin{enumerate}[label=\numberballred]
    \item \textbf{Grupos y tema.} En pequeños grupos, elegid un tema distinto para cada grupo.
    \item \textbf{Investigación.} Investigad las circunstancias históricas que impulsaron esos avances (soluciones de ingeniería, nuevos materiales, descubrimientos previos) y las posibles resistencias sociales o políticas.
    \item \textbf{Línea del tiempo.} Elegid un modelo y situad en ella las etapas (investigaciones, descubrimientos, teorías, leyes, publicaciones, controversias), sus protagonistas y los sucesos históricos de cada momento.
    \item \textbf{Presentación.} Preparad una presentación sobre los científicos y científicas que contribuyeron al tema, con sus datos biográficos más importantes e ilustraciones (citando su procedencia).
    \item \textbf{Comunicación.} Presentad la línea del tiempo en clase, comparad hechos ocurridos en la misma época en los distintos temas y publicad los trabajos en la revista del departamento.
\end{enumerate}

%========================================================================================
\section{Unidad 2. La historia del átomo en una línea del tiempo}
%========================================================================================
% (sacado del final de UD_2_el_atomo_tp.tex)

En 1869, el químico ruso \textbf{Dmitri Mendeléiev} ordenó los 63 elementos conocidos por masa atómica creciente y se dio cuenta de que las propiedades se repetían periódicamente. Colocó juntos los elementos parecidos y, cuando un elemento no encajaba, cambió el orden o \textbf{dejó un hueco} para un elemento aún desconocido, cuyas propiedades predijo. Así ocurrió con el galio (descubierto en 1875) y el germanio (1886). Entre 1894 y 1898, \textbf{William Ramsay} descubrió los gases nobles, que no estaban previstos, y los colocó en una nueva columna: el actual grupo 18.

\begin{marker}
    \textbf{Reto:} diseñad en grupos una \textbf{línea del tiempo} que recoja los principales descubrimientos que llevaron a los distintos modelos atómicos, los relacionados con la radiactividad y las fechas de descubrimiento de algunos elementos químicos.
\end{marker}

\begin{enumerate}[label=\numberballred]
    \item \textbf{Investigación.} Repartid el trabajo en grupos de dos o tres personas sin que se repitan los temas. Anotad las fechas que aparecen en la unidad y ampliadlas con modelos más modernos, la radiactividad y sus aplicaciones, y el descubrimiento de elementos. Identificad a los científicos y científicas protagonistas.
    \item \textbf{Fichas.} Diseñad una ficha común para recoger de forma homogénea cada acontecimiento (fecha, protagonista, descubrimiento, importancia, imagen). Añadid algún hecho histórico que sirva de referencia de cada época.
    \item \textbf{Construcción.} Elegid un formato (mural o herramienta digital) adaptado al lugar donde se expondrá y acompañad la información con ilustraciones. Haced un reportaje fotográfico o en vídeo del proceso.
    \item \textbf{Comunicación.} Exponed la línea del tiempo en el aula o en el centro, explicadla a compañeros de otros cursos y publicadla en la revista digital, el blog de la asignatura o la web del departamento.
\end{enumerate}

%========================================================================================
\section{Unidad 3. Materiales para la ingeniería y el deporte}
%========================================================================================
% (sacado del final de UD3_enlace_quimico.tex)

Los deportistas baten sus propias marcas una y otra vez. En buena parte se debe a los \textbf{nuevos materiales}: más flexibles y ligeros, pero también más fuertes y resistentes. La madera, el hierro o el cuero han ido dejando paso a compuestos químicos de alta tecnología. Para saltar más alto, nadar más rápido, bucear más hondo o correr más deprisa hacen falta cuerdas más ligeras, trajes protectores, zapatillas especiales, raquetas, pelotas o palos de golf más ligeros, fuertes y elásticos. Muchas veces, el resultado de un atleta depende de la calidad tecnológica de su equipamiento: la velocidad es, en definitiva, una cuestión de física y de química.

\begin{marker}
    \textbf{Reto:} investigad los materiales (elementos y compuestos) que se usan en la ingeniería y la construcción, y los nuevos materiales de las equipaciones e instrumentos deportivos. Preparad unas \textbf{fichas identificativas} y una \textbf{exposición física o virtual}. Podéis hacerlo en colaboración con el departamento de Educación Física.
\end{marker}

\begin{enumerate}[label=\numberballred]
    \item \textbf{Investigación.} Buscad información sobre los materiales que se usan en ingeniería, construcción y edificación. No olvidéis las nuevas aleaciones metálicas ni las formas alotrópicas del carbono (fibra de carbono, grafeno, nanotubos).
    \item \textbf{Fichas.} Diseñad una ficha común que recoja las características de los materiales de las equipaciones deportivas y de los instrumentos de cada deporte (tipo de sustancia, enlace, propiedades, ventajas). Si es posible, comparadlos con los que se usaban hace unas décadas.
    \item \textbf{Análisis.} Elegid un deporte y comparad sus marcas de hace 20 o 30 años con las actuales. ¿Puede estar relacionada la mejora con el uso de nuevos materiales?
    \item \textbf{Elaboración.} Haced una lista de materiales de ingeniería y construcción y sus aplicaciones, acompañada de fotografías o ilustraciones.
    \item \textbf{Comunicación.} Preparad en clase una presentación, a modo de reportaje, sobre la evolución de los materiales deportivos y su relación con la evolución de las marcas, y exponedla en el centro o en la web del departamento.
\end{enumerate}

%========================================================================================
\section[Unidad 4. Compuestos de carbono y plásticos]{Unidad 4. Compuestos de carbono de interés industrial y biológico. Reciclaje de plásticos}
%========================================================================================
% (sacado del final de UD4_organica.tex)

Se conocen más de diez millones de compuestos del carbono, y gran parte de la industria química se dedica a fabricar plásticos. Pero sus residuos son un problema: los envases de un solo uso son uno de los elementos más contaminantes. En 2022, la Comisión Europea propuso nuevas normas porque cada europeo genera unos \qty{180}{\kg} de residuos de envases al año, y sin medidas los residuos de envases de plástico podrían aumentar un \qty{46}{\percent} para 2030. Entre las medidas están prohibir algunos envases de un solo uso (para comer y beber en los restaurantes, para frutas y verduras o los envases en miniatura de los hoteles), exigir que todos los envases sean reciclables y crear sistemas de devolución de botellas y latas. El objetivo es \textbf{prevenir} la generación de envases, \textbf{impulsar} un reciclado de calidad y \textbf{reducir} el consumo de recursos naturales.

\begin{marker}
    \textbf{Reto:} investigar las propiedades y aplicaciones de un compuesto de carbono cercano y analizar el consumo de envases no reutilizables en el centro, con una campaña de concienciación (eslóganes y póster).
\end{marker}

\begin{enumerate}[label=\numberballred]
    \item \textbf{Grupos y compuesto.} En pequeños grupos, elegid un compuesto de carbono (etanol, gas natural, neopreno, aceites, glucosa, PVC, vitaminas, jabón, caucho, percloroetileno\ldots), procurando que no se repitan entre los grupos.
    \item \textbf{Investigación.} Buscad si es natural o artificial, sus propiedades físicas y químicas, sus aplicaciones (domésticas, clínicas, industriales), los problemas ambientales que puede causar y las normas que regulan su uso.
    \item \textbf{Ficha común.} Diseñad entre todos un modelo de ficha: nombre, fórmula, masa molecular, puntos de fusión y ebullición, aplicaciones, importancia biológica y problemas para la salud o el medio ambiente. Elaborad con las fichas un póster mural o una infografía.
    \item \textbf{Los plásticos del centro.} Analizad qué envases de plástico se generan en el centro, comprobad si se depositan en el contenedor adecuado y redactad un informe sobre cómo reducirlos o sustituirlos por envases reutilizables.
    \item \textbf{Campaña.} Diseñad pósteres con eslóganes sobre el uso excesivo de plásticos y su efecto en el medio ambiente, y publicad los resultados en la revista digital del departamento.
    \item \textbf{Evaluación.} Al cabo de un tiempo, comprobad si han disminuido los residuos plásticos del centro (podéis hacer también una encuesta al alumnado de otros grupos).
\end{enumerate}

\begin{exercise}
    El PVC (cloruro de polivinilo) se obtiene por polimerización del cloruro de vinilo, \ch{CH2=CHCl}. a) Calcula la masa molecular del monómero y su porcentaje de cloro. b) ¿Qué tipo de plástico es el PVC según su comportamiento con el calor? c) ¿Por qué es un problema que los plásticos sean tan estables?
    \begin{solu}
        a) \ch{C2H3Cl}: \(M = 24 + 3 + \num{35.5} = \textcolor{red}{\num{62.5}}\); \% Cl \(= \dfrac{\num{35.5}}{\num{62.5}} \cdot 100 = \textcolor{red}{\qty{56.8}{\percent}}\).
        b) Es un \textcolor{red}{termoplástico}: se reblandece al calentarlo y puede moldearse y reciclarse.
        c) Porque \textcolor{red}{tardan cientos de años en degradarse}: se acumulan en vertederos, ríos y océanos, y se fragmentan en microplásticos que entran en las cadenas alimentarias.
    \end{solu}
\end{exercise}

%========================================================================================
\section{Unidad 5. La atmósfera, un sistema material que proteger}
%========================================================================================
% (sacado del final de UD5_Sistema_materiales.tex)

La contaminación atmosférica se debe a la emisión de gases y partículas procedentes de la actividad humana y de fuentes naturales, como los volcanes. Algunos contaminantes que afectan a nuestra salud y a los ecosistemas son el dióxido de azufre (\ch{SO2}), el trióxido de azufre (\ch{SO3}), los óxidos de nitrógeno (\ch{NO2}, \ch{N2O}\ldots), el monóxido de carbono (\ch{CO}), el dióxido de carbono (\ch{CO2}) y el ozono que se forma a baja altura. También contaminan las \textbf{partículas en suspensión}: polvo, metales y compuestos orgánicos e inorgánicos. El clima influye en cómo se dispersan y transportan estos contaminantes.

\begin{marker}
    \textbf{Reto:} evaluad la calidad del aire de vuestro entorno en distintas épocas del año y redactad un \textbf{decálogo de buenas prácticas} para reducir los contaminantes atmosféricos.
\end{marker}

\begin{enumerate}[label=\numberballred]
    \item \textbf{La capa inicial de la atmósfera.} El \qty{95}{\percent} de la masa del aire (unos \qty{5.1e18}{\kg}) está en los primeros \qty{11}{\km}. Investigad en pequeños grupos qué gases la componen y en qué proporción.
    \item \textbf{Datos de calidad del aire.} Con alguna aplicación o web de calidad del aire, anotad durante varias semanas, en distintas épocas del año, la calidad del aire de vuestra localidad y el contaminante principal. ¿Coinciden esos gases con los de la composición del aire?
    \item \textbf{Código de colores.} Analizad qué significan los colores (verde, amarillo, rojo\ldots) con que se indica la calidad del aire.
    \item \textbf{Zonas de bajas emisiones.} Averiguad si en vuestra localidad hay zonas de bajas emisiones (ZBE), qué vehículos pueden entrar y por qué. Situadlas en un plano y comprobad si en ellas la calidad del aire es mejor.
    \item \textbf{Comunicación.} Representad los datos en tablas y gráficas y publicad el análisis en el centro o en la revista digital del departamento.
    \item \textbf{Decálogo.} En grupos, redactad buenas prácticas individuales y colectivas para reducir la contaminación del aire. Ponedlas en común, elaborad un único decálogo de la clase y exponedlo en el aula.
\end{enumerate}

%========================================================================================
\section{Unidad 6. Las combustiones en nuestra vida diaria}
%========================================================================================
% (sacado del final de UD6_reacciones_quimicas.tex)

Piensa en cualquier producto que consumes: hay que fabricarlo, envasarlo, almacenarlo, transportarlo hasta la tienda, anunciarlo y, al final, deshacerse de él. Todo ello genera residuos y consume mucha energía, buena parte procedente de combustiones. Comprar productos locales reduce las distancias del transporte y, con ello, el consumo de energía y la emisión de gases de efecto invernadero. Además, el comercio de proximidad ayuda a que las ciudades sean más sostenibles. Aunque las grandes decisiones corresponden a los gobiernos, los ODS solo se alcanzarán con acciones individuales comprometidas.

\begin{marker}
    \textbf{Reto:} estudiad el impacto medioambiental de las combustiones que usamos para obtener energía en casa y en la industria, y realizad una \textbf{campaña informativa} sobre los efectos de quemar combustibles fósiles, entre ellos la lluvia ácida y el incremento del efecto invernadero.
\end{marker}

\begin{enumerate}[label=\numberballred]
    \item \textbf{Investigación.} Averiguad cómo se usan las combustiones para producir electricidad y cómo se emplea esa energía en el transporte, en los hogares y en los lugares de trabajo, estudio y ocio.
    \item \textbf{Transporte de mercancías.} Buscad datos sobre las combustiones necesarias para transportar personas y mercancías, y para fabricar materias primas y bienes de consumo a veces muy lejos del lugar donde se usan.
    \item \textbf{Análisis.} Comparad el doble impacto (producción y transporte) de los productos lejanos con el de los productos de proximidad.
    \item \textbf{Póster 1.} Mostrad que somos responsables de la contaminación que producimos, tanto con el transporte como con la energía doméstica, incluyendo la que generan los productos que vienen de lejos.
    \item \textbf{Póster 2.} Explicad cómo se produce la electricidad en una central térmica y qué combustible usa. Recordad que esa energía alimenta todos los aparatos eléctricos, incluidos los vehículos eléctricos.
    \item \textbf{Decálogo.} Redactad buenas prácticas individuales y colectivas para reducir el impacto de los combustibles fósiles y publicadlas en la revista del centro o del departamento.
    \item \textbf{Comunicación.} Exponed la campaña a toda la comunidad educativa, con los pósteres como apoyo.
\end{enumerate}

%========================================================================================
\section{Unidad 7. El movimiento en el deporte}
%========================================================================================
% (sacado del final de UD7_movimientos_rectilineos.tex)

En el deporte hay que medir distancias en los campos de juego, en los saltos y en los lanzamientos, y tiempos en las carreras a pie o de patinaje. Muchas medidas se hacen fuera de la pista, en laboratorios que prueban los materiales. En la Fórmula 1 se miden longitudes, tiempos y masas para cumplir el reglamento, y en la natación artística se controlan la temperatura y la profundidad del agua. Para homologar un récord, las medidas deben hacerse con instrumentos certificados por un laboratorio reconocido.

\begin{marker}
    \textbf{Reto:} investigad cómo se miden las magnitudes del movimiento (sistema de referencia, espacio, tiempo, velocidad, aceleración) en distintos deportes y elaborad una \textbf{infografía} o un \textbf{pódcast} con los resultados.
\end{marker}

\begin{enumerate}[label=\numberballred]
    \item \textbf{Análisis.} En grupos de hasta cinco personas, haced una lista de deportes en los que lo más importante sea la velocidad y otra en los que lo sea la resistencia. Cada miembro elige una disciplina diferente de la primera lista.
    \item \textbf{Investigación.} Buscad cómo se establece el sistema de referencia (salida y llegada), cómo se mide el espacio recorrido y el tiempo, y qué se necesita para que una marca sea homologable.
    \item \textbf{Récords.} Consultad las marcas mundiales actuales de vuestra disciplina y cómo han evolucionado en las dos últimas décadas.
    \item \textbf{Nuestra carrera.} Con el profesorado de Educación Física, organizad una carrera de \qty{100}{\m} con toda la clase: calculad la velocidad media de cada participante y analizad las fases de la carrera (salida, aceleración, velocidad máxima y resistencia), por ejemplo midiendo los tiempos de paso cada \qty{20}{\m}.
    \item \textbf{Ficha común.} Diseñad entre toda la clase una ficha con la misma estructura para recoger la información y cumplimentadla.
    \item \textbf{Infografía o pódcast.} Elaborad una infografía o un pódcast que resuma vuestra investigación, y un informe personal con los resultados de la carrera.
    \item \textbf{Comunicación.} Presentad vuestro trabajo al resto de la clase y debatid sobre los métodos de medida en los deportes de velocidad.
\end{enumerate}

%========================================================================================
\section{Unidad 8. El rozamiento y la acción-reacción en la vida diaria}
%========================================================================================
% (sacado del final de UD8_fuerzas.tex)

En la final del Mundial de natación de Roma, el 30 de julio de 2009, el brasileño César Cielo tocó la pared en \qty{46.91}{\s} en los \qty{100}{\m} libres, ayudado por un bañador de poliuretano que reducía el rozamiento con el agua y lo hacía flotar mejor. Esos bañadores se prohibieron en 2010 y, durante años, ningún nadador con bañador de tela consiguió acercarse a aquella marca. En agosto de 2022, el rumano David Popovici la batió con \qty{46.86}{\s}: rotaba el tronco y clavaba las manos en el agua, ofreciendo muy poca superficie al rozamiento.

\begin{marker}
    \textbf{Reto:} identificad e interpretad las fuerzas de rozamiento y las fuerzas de acción y reacción en situaciones cotidianas y deportivas, y preparad una \textbf{presentación}, una \textbf{infografía} o un \textbf{mural}.
\end{marker}

\begin{enumerate}[label=\numberballred]
    \item \textbf{Grupos.} Formad grupos de hasta cinco personas.
    \item \textbf{Situaciones.} Buscad situaciones de la vida diaria y del deporte en las que aparezcan fuerzas de rozamiento y de acción y reacción: caminar sobre distintos terrenos, borrar la pizarra, arrastrar un objeto, saltar desde una barca o un monopatín, nadar\ldots
    \item \textbf{Análisis.} Identificad en cada caso las fuerzas de rozamiento y las parejas de acción y reacción, y dibujadlas.
    \item \textbf{Los bañadores.} ¿Por qué se conseguían mejores marcas con los bañadores de poliuretano que con los de tela? ¿Creéis que estuvo justificada su prohibición?
    \item \textbf{Técnica.} Analizad cómo nadaba Popovici para ser más rápido y eficaz.
    \item \textbf{Puesta en común.} Haced una lista con todas las situaciones estudiadas para que no se repitan entre grupos, y comprobad cuántas se aprovechan para facilitar el movimiento, sobre todo en el deporte.
    \item \textbf{Elaboración.} Buscad fotografías o vídeos (citando su autoría) y preparad la presentación, la infografía o el mural con al menos cuatro situaciones.
    \item \textbf{Comunicación.} Presentad vuestro trabajo a la clase o difundidlo con los medios del centro.
\end{enumerate}

%========================================================================================
\section[Unidad 9. El universo y la gravedad]{Unidad 9. Conocimiento del universo y la fuerza gravitatoria}
%========================================================================================
% (sacado del final de UD9_movimiento_circular.tex)

La gravedad influye en el desarrollo de las plantas en la Tierra: las raíces crecen hacia abajo y los tallos hacia arriba. En el experimento \emph{Plant Gravity Sensing}, realizado en la Estación Espacial Internacional, se cultivaron semillas de \emph{Arabidopsis} (una planta muy usada en investigación) en microgravedad. Sus raíces crecieron enredadas en la malla de cultivo, de forma muy distinta a las cultivadas en la Tierra. Entender cómo «sienten» las plantas la gravedad podría ayudar a obtener nuevas variedades de cultivos.

Investigaciones como esta son posibles gracias a siglos de avances en el conocimiento del universo y de la gravedad, desde Copérnico y Kepler hasta los satélites y las estaciones espaciales.

\begin{marker}
    \textbf{Reto:} construid una \textbf{línea del tiempo} con la evolución del conocimiento del universo y de la fuerza gravitatoria, de los métodos de observación del universo y de la navegación espacial, junto con sus protagonistas más destacados.
\end{marker}

\begin{enumerate}[label=\numberballred]
    \item \textbf{Grupos y tema.} En grupos de hasta cinco personas, elegid uno de estos aspectos: evolución del conocimiento del universo; evolución de las teorías sobre la gravedad; métodos de observación del universo; la navegación espacial\ldots{} Procurad que no haya más de dos grupos con el mismo aspecto.
    \item \textbf{Investigación.} Ampliad la información sobre los experimentos que se hacen en la Estación Espacial Internacional y cómo contribuyen al avance de la ciencia (podéis consultar la web de la ESA).
    \item \textbf{Ficha común.} Diseñad entre todos un modelo de ficha para recoger la información de cada hecho científico en su contexto histórico: fecha, protagonistas, aportación, contexto y fuentes.
    \item \textbf{Elaboración.} Decidid dónde vais a exponer la línea del tiempo y elegid las dimensiones adecuadas. Buscad un modelo o diseñad el vuestro e incluid ilustraciones de los hechos y de sus protagonistas, citando siempre su procedencia.
    \item \textbf{Comunicación.} Exponed la línea del tiempo en clase; todos los miembros del grupo deben participar. Después, elegid el hecho que os parezca más importante y explicad por qué.
    \item \textbf{Debate y difusión.} Debatid sobre la trascendencia de los hechos más destacados y explicad a alumnado de otros cursos qué significa vuestra línea del tiempo y cómo la habéis hecho.
\end{enumerate}

%========================================================================================
\section[Unidad 10. Presión atmosférica y mapas del tiempo]{Unidad 10. Presión atmosférica, fenómenos meteorológicos y mapas del tiempo}
%========================================================================================
% (sacado del final de UD10_fluidos.tex)

En 1648, \textbf{Blaise Pascal} quiso comprobar si la columna de mercurio del tubo de Torricelli se sostenía por el peso del aire. Si era así, en lo alto de una montaña, con menos aire encima, la columna tenía que ser más corta. Le pidió a su cuñado, Florin Périer, que hiciera la prueba en el Puy-de-Dôme, una montaña de Francia. Périer llenó dos tubos iguales y comprobó al pie de la montaña que en ambos el mercurio quedaba a la misma altura, unos \qty{71}{\cm}. Dejó uno de ellos allí, vigilado, y subió con el otro casi \qty{1000}{\m}: en la cima, el mercurio solo alcanzaba unos \qty{63}{\cm}. Al bajar, la columna volvía a crecer. El aire pesa, y la presión que ejerce disminuye con la altitud.

Hoy la presión atmosférica es una de las magnitudes básicas de la meteorología. ¿Sabrías distinguir una borrasca de un anticiclón en un mapa del tiempo? ¿Y un frente frío de uno cálido?

\begin{marker}
    \textbf{Reto:} aplicar lo que sabéis de la presión atmosférica para describir los fenómenos meteorológicos e interpretar mapas del tiempo, reconociendo los términos y los símbolos propios de la meteorología. Al final, debatiréis vuestras conclusiones y redactaréis un informe.
\end{marker}

\begin{enumerate}[label=\numberballred]
    \item \textbf{Grupos.} Formad grupos de trabajo de un máximo de cinco personas.
    \item \textbf{Análisis.} ¿Qué magnitud estaba midiendo Pascal con el tubo de mercurio? ¿Qué nombre recibe ese instrumento? ¿Qué diferencia de presión, en \unit{\mmHg} y en \unit{\hecto\pascal}, midió Périer entre el pie y la cima de la montaña?
    \item \textbf{Investigación.} Buscad los términos y los símbolos más usados en meteorología y en los mapas del tiempo (milibar, hectopascal, isobara, anticiclón, borrasca, frente frío, frente cálido, superficie frontal\ldots), y la relación entre la presión en las capas altas y bajas de la atmósfera y el tiempo.
    \item \textbf{Elaboración.} Redactad una ficha para cada símbolo de los mapas del tiempo y otra para cada término que aparece con frecuencia en las noticias meteorológicas.
    \item \textbf{Interpretación.} Descargad mapas de isobaras de España durante una semana (por ejemplo, de la web de AEMET), interpretadlos y describid el tiempo que predicen. Comparad también mapas de vuestra comunidad autónoma en distintas estaciones del año.
    \item \textbf{Comunicación.} Preparad una presentación, una exposición oral y un póster con las fichas y los mapas que habéis interpretado, y publicad vuestro informe y vuestras conclusiones en la revista del departamento o en el aula virtual.
\end{enumerate}

\begin{exercise}
    Responde a las preguntas del paso 2 de la situación de aprendizaje.
    \begin{solu}
        Pascal medía la \textcolor{red}{presión atmosférica}, y el instrumento es un \textcolor{red}{barómetro}. La diferencia de presión es la de una columna de unos \(71 - 63 = \qty{8}{\cm} = \qty{80}{\mmHg}\):
        \(\Delta p = 13600 \cdot \num{9.8} \cdot \num{0.080} = \qty{1.07e4}{\pascal} \approx \textcolor{red}{\qty{107}{\hecto\pascal}}\).
    \end{solu}
\end{exercise}

%========================================================================================
\subsection[Proyecto: las fuerzas de la naturaleza]{Proyecto: las fuerzas de la naturaleza en nuestro entorno}
%========================================================================================

¿Qué pasaría en nuestra vida diaria si alguna de las fuerzas de la naturaleza, o de las leyes que las rigen, dejara de cumplirse? En este proyecto, que cierra el bloque de las fuerzas, vais a escribir un \textbf{relato breve de ciencia ficción} sobre un mundo en el que una ley física no funciona.

\begin{enumerate}[label=\numberballred]
    \item \textbf{Investigación.} Haced una lista de las leyes que hemos estudiado en las unidades de este bloque: las tres leyes de Newton, la fuerza de rozamiento, la fuerza elástica, la fuerza centrípeta, la ley de gravitación universal, la presión hidrostática, el principio de Pascal, el principio de Arquímedes\ldots{} Formad grupos de tres o cuatro personas y repartid las leyes entre los grupos.
    \item \textbf{Lluvia de ideas.} Cada grupo imagina qué sucedería si no se cumpliera su ley: ¿a qué hechos o fenómenos daría lugar? Tomad nota de todas las ideas y analizadlas una a una.
    \item \textbf{Elaboración.} Preparad un pequeño índice o esquema para ordenar las ideas y repartir el trabajo. Después, redactad el relato (unas quince páginas como máximo). Buscad información en fuentes fiables y citadlas; usad imágenes propias o con licencia libre.
    \item \textbf{Comunicación.} Presentad el relato en clase justificando la física que hay detrás. Podéis organizar un concurso en el que se valoren la originalidad y el rigor científico de cada relato, y publicarlos en el aula virtual o en la revista del departamento.
    \item \textbf{Reflexión.} Comentad en el grupo qué dificultades habéis encontrado y qué habéis aprendido. Relacionad el proyecto con el \textbf{ODS 5} (\emph{Igualdad de género}), repartiendo las tareas y el protagonismo de forma equitativa, y con el \textbf{ODS 11} (\emph{Ciudades y comunidades sostenibles}): ¿cómo cambiarían el transporte o la seguridad vial en vuestro mundo imaginario?
\end{enumerate}

%========================================================================================
\section[Unidad 11. La energía en las ferias]{Unidad 11. La conservación de la energía en ferias y parques de atracciones}
%========================================================================================
% (sacado del final de UD11_energia.tex)

Las montañas rusas nacieron en San Petersburgo, donde la gente se deslizaba en trineos por grandes toboganes de madera cubiertos de nieve; en Francia, donde se copió la idea con vagonetas sobre raíles, se la llamó «montaña rusa». A finales del siglo \textsc{xix} se construyeron las primeras con bucles, y a comienzos del siglo \textsc{xx} se añadieron ruedas por debajo de la vía, que permiten curvas muy pronunciadas a gran velocidad.

En casi todas, una cadena motorizada sube el tren hasta la primera cima y lo suelta: a partir de ahí, el recorrido depende solo de la \textbf{conservación de la energía}. La energía potencial de la cima se transforma en cinética en las bajadas y vuelve a ser potencial en las subidas, mientras el rozamiento con la vía y el aire va disipando poco a poco energía. Se podría decir que solo pagamos por la subida; el resto lo hace la física.

\begin{marker}
    \textbf{Reto:} identificar las transformaciones de la energía mecánica que tienen lugar en las atracciones de los parques de ocio y comprobar que su funcionamiento se basa en el principio de conservación de la energía. Al final, debatiréis vuestras conclusiones y redactaréis un informe.
\end{marker}

\begin{enumerate}[label=\numberballred]
    \item \textbf{Grupos.} Formad grupos de trabajo de un máximo de cuatro personas.
    \item \textbf{Investigación.} Elegid varias atracciones (buscando información en internet) y haced una ficha para cada una: nombre, tipo (caída libre, bucle, barca pendular, montaña rusa, rampa con final en el agua, coches de choque\ldots) y qué energía se usa para ponerla en marcha.
    \item \textbf{Hipótesis.} Proponed cómo funciona cada atracción y qué transformaciones de energía tienen lugar. Ayudaos de dibujos, esquemas, fotografías o vídeos.
    \item \textbf{Elaboración.} Justificad vuestras hipótesis con el principio de conservación de la energía: energía inicial, transformaciones durante el recorrido, pérdidas por rozamiento o ruido, sistema de frenado\ldots{} Razonad por qué es tan importante la altura del punto de partida.
    \item \textbf{Comunicación.} Preparad un informe y una presentación con la secuencia ordenada de transformaciones de cada atracción. Debatid con los demás grupos, recoged las conclusiones y publicadlas en la revista del departamento o en el aula virtual.
\end{enumerate}

\begin{exercise}
    En una montaña rusa, el tren (con sus pasajeros, \qty{3000}{\kg}) parte casi del reposo desde la primera cima, a \qty{40}{\m} de altura.
    \begin{enumerate}[label=\alph*)]
        \item ¿Qué energía tiene que aportar la cadena para subir el tren hasta la cima?
        \item Sin rozamiento, ¿con qué velocidad pasaría por el punto más bajo, a nivel del suelo?
        \item ¿Podría la vía tener una segunda cima de \qty{45}{\m}? ¿Por qué la primera cima siempre es la más alta?
        \item En realidad, el tren pasa por el punto más bajo a \qty{25}{\m\per\s}. ¿Cuánta energía se ha transformado en energía térmica?
    \end{enumerate}
    \begin{solu}
        a) \(E_p = 3000 \cdot \num{9.8} \cdot 40 = \textcolor{red}{\qty{1.18e6}{\joule}}\).
        b) \(v = \sqrt{2 \cdot g \cdot h} = \sqrt{2 \cdot \num{9.8} \cdot 40} = \textcolor{red}{\qty{28}{\m\per\s}}\) (unos \qty{100}{\km\per\hour}).
        c) \textcolor{red}{No}: el tren no tiene energía suficiente para subir más alto que el punto de partida (y menos aún con rozamiento). Por eso la primera cima es siempre la más alta.
        d) \(E_c = \frac{1}{2} \cdot 3000 \cdot 25^2 = \qty{9.4e5}{\joule}\); la energía térmica es \(\num{1.18e6} - \num{9.4e5} = \textcolor{red}{\qty{2.4e5}{\joule}}\).
    \end{solu}
\end{exercise}

%========================================================================================
\section[Unidad 12. La energía que consumimos]{Unidad 12. Estimación de la energía que consumimos}
%========================================================================================
% (sacado del final de UD12_calor.tex)

El Sol emite cada segundo una energía de unos \qty{1e20}{\kilo\watt\hour}. Solo una parte diminuta llega a la Tierra, y de ella algo menos de la mitad alcanza la superficie: unos \qty{7e17}{\kilo\watt\hour} al año. La humanidad consume cada año unos \qty{1.7e14}{\kilo\watt\hour} de energía comercial (la que se compra y se vende), a la que habría que sumar la energía «no comercial», como la leña que se recoge para cocinar. Aun así, la energía que consumimos es solo una parte diminuta de la que el Sol deja cada año sobre la superficie terrestre.

Ponte en situación y haz las cuentas: ¿cuál es la primera conclusión que sacarías?

\begin{marker}
    \textbf{Reto:} estimar la energía que consumimos en casa o en el instituto y proponer medidas de ahorro energético, recogidas en un \textbf{decálogo de buenas prácticas} y acompañadas de \textbf{eslóganes} que animen a usar la energía de forma racional. Necesitaréis facturas de la electricidad y de los combustibles (gas natural, gasóleo, bombonas de butano) de casa o del instituto.
\end{marker}

\begin{enumerate}[label=\numberballred]
    \item \textbf{Análisis de la factura de la luz.} Localizad el consumo del período facturado (en \unit{\kilo\watt\hour}) y el importe total. Observad que la factura incluye otros conceptos además de la energía consumida (potencia contratada, impuestos, alquiler del contador\ldots). Calculad el precio medio del \unit{\kilo\watt\hour} y el consumo medio diario. Buscad el consumo de todo el año.
    \item \textbf{Combustibles.} Haced lo mismo con la factura del gas natural, del gasóleo o del butano que se use para cocinar, calentar agua o la calefacción. En el gas natural, la factura indica el consumo en \unit{\kilo\watt\hour}: cada \unit{\cubic\m} de gas equivale a unos \qty{11.7}{\kilo\watt\hour}.
    \item \textbf{Investigación.} Haced una lista de los aparatos eléctricos que más contribuyen al consumo, con su potencia y su tiempo de uso diario, y estimad su consumo y su coste mensual. Averiguad cuántas horas a la semana está encendida la calefacción del instituto y estimad el gasto de gas en casa (cocina, baño, calefacción).
    \item \textbf{Renovables.} Averiguad si el instituto o vuestra casa tienen instalaciones de energías renovables (paneles fotovoltaicos, colectores solares\ldots) y si la energía obtenida se usa en el propio edificio o se vende a la red.
    \item \textbf{Elaboración y comunicación.} En grupos pequeños, elaborad una lista de medidas viables para reducir el consumo en casa y en el instituto. Ponedlas en común, redactad el decálogo y los eslóganes y organizad una campaña de sensibilización en el centro. Al cabo de un mes, comprobad si ha bajado el consumo.
\end{enumerate}

\begin{exercise}
    Comprueba las cuentas del texto y responde:
    \begin{enumerate}[label=\alph*)]
        \item ¿Qué fracción de la energía solar que llega cada año a la superficie terrestre consume la humanidad?
        \item ¿Qué conclusión sacas?
    \end{enumerate}
    \begin{solu}
        a) \(\dfrac{\num{1.7e14}}{\num{7e17}} \approx \num{2.4e-4}\), es decir, unas \textcolor{red}{dos o tres diezmilésimas} partes (aproximadamente, la energía que llega a la superficie en solo dos horas). Aunque se sume la energía no comercial, la proporción sigue siendo muy pequeña.
        b) La energía solar que recibe la Tierra es \textcolor{red}{miles de veces mayor} que la que consume la humanidad: en principio, bastaría para cubrir todas nuestras necesidades si fuéramos capaces de aprovecharla.
    \end{solu}
\end{exercise}

\begin{exercise}
    En una factura bimestral de la luz aparece un consumo de \qty{420}{\kilo\watt\hour} y un importe total de 126 euros.
    \begin{enumerate}[label=\alph*)]
        \item Calcula el precio medio de cada \unit{\kilo\watt\hour} facturado y el consumo medio diario (toma 60 días).
        \item En casa hay un horno eléctrico de \qty{2000}{\watt} que se usa \qty{30}{\minute} al día. ¿Qué parte del consumo diario supone?
        \item La misma vivienda consume en ese período \qty{45}{\cubic\m} de gas natural. ¿A cuántos \unit{\kilo\watt\hour} equivale?
    \end{enumerate}
    \begin{solu}
        a) Precio medio: \(\dfrac{126}{420} = \textcolor{red}{\num{0.30}}\) euros por \unit{\kilo\watt\hour}. Consumo diario: \(\dfrac{420}{60} = \textcolor{red}{\qty{7}{\kilo\watt\hour}}\).
        b) \(E = P \cdot t = 2 \cdot \num{0.5} = \qty{1}{\kilo\watt\hour}\) diario: \(\dfrac{1}{7} \approx \textcolor{red}{\qty{14}{\percent}}\) del consumo.
        c) \(45 \cdot \num{11.7} \approx \textcolor{red}{\qty{527}{\kilo\watt\hour}}\).
    \end{solu}
\end{exercise}

%========================================================================================
\section[Unidad 13. STOP a la contaminación acústica]{Unidad 13. STOP a la contaminación acústica}
%========================================================================================
% (sacado del final de UD13_luz_sonido.tex)

Un estudio del Instituto de Salud Global de Barcelona (ISGlobal), publicado en 2022, concluyó que el alumnado de escuelas con más ruido de tráfico muestra un desarrollo cognitivo más lento: el ruido perjudica la \textbf{memoria de trabajo} (la que nos permite mantener y manejar información durante un rato) y la \textbf{capacidad de atención} (la que nos permite concentrarnos en una tarea). En el estudio, cada aumento de \qty{5}{\decibel} del ruido exterior se relacionó con una disminución de la capacidad de atención.

\begin{marker}
    \textbf{Reto:} concienciar a toda la comunidad educativa de que el nivel de intensidad sonora en el centro no debería superar los \qty{35}{\decibel}. Para ello mediréis el nivel sonoro en distintos lugares del centro (aulas, biblioteca, pasillos, escaleras, patio, gimnasio, laboratorios, talleres\ldots) y en distintos momentos de la jornada (entradas y salidas, cambios de clase, recreo, durante las clases\ldots), y elaboraréis un \textbf{decálogo de buenas prácticas} ilustrado con imágenes.
\end{marker}

\textbf{Análisis e investigación}
\begin{enumerate}[label=\numberballred]
    \item Elegid entre toda la clase el sonómetro que vais a usar (analógico, digital o una aplicación del móvil). Todos los grupos deben usar el mismo, para que las medidas sean comparables.
    \item Elegid los lugares y los momentos concretos de cada medida, de forma que sean comprobables y que todos los grupos midan en las mismas condiciones.
    \item Diseñad una plantilla de toma de datos en la que figuren el lugar, el día, la hora y el nivel medido.
\end{enumerate}

\textbf{Elaboración y comunicación}
\begin{enumerate}[label=\numberballred, resume]
    \item Haced un reportaje fotográfico del desarrollo de la investigación.
    \item Con todos los datos, elaborad \textbf{mapas de ruido} del centro (sobre un plano) para los distintos momentos de la jornada. Destacad con colores las zonas y los momentos en los que el nivel sea igual o superior a \qty{35}{\decibel}.
    \item Exponed los mapas en un lugar visible para toda la comunidad educativa.
    \item Elaborad eslóganes y un decálogo de buenas prácticas con medidas \textbf{activas} (que reducen el ruido en su origen: hablar más bajo, no arrastrar las sillas, poner tacos de goma en las patas\ldots) y \textbf{pasivas} (que impiden que el ruido se propague: cerrar puertas y ventanas, paneles absorbentes, cortinas\ldots).
    \item Publicad el decálogo en las aulas y en los tablones del centro. Pasado un tiempo, repetid las medidas para comprobar si la campaña ha tenido efecto.
\end{enumerate}

\begin{exercise}
    En un pasillo se miden \qty{75}{\decibel} durante un cambio de clase y \qty{45}{\decibel} durante las clases. ¿Cuántas veces mayor es la intensidad física del sonido en el cambio de clase?
    \begin{solu}
        La diferencia es de \(75 - 45 = \qty{30}{\decibel}\): tres veces \qty{10}{\decibel}. La intensidad es \(10 \cdot 10 \cdot 10 = \textcolor{red}{1000}\) veces mayor.
    \end{solu}
\end{exercise}

\begin{exercise}
    Clasifica estas medidas en activas o pasivas: a) poner gomas en las patas de las sillas; b) instalar ventanas de doble cristal; c) hablar en voz baja en la biblioteca; d) colocar paneles de corcho en las paredes del aula de música.
    \begin{solu}
        Activas (reducen el ruido en el foco): \textcolor{red}{a y c}. Pasivas (dificultan su propagación o lo absorben): \textcolor{red}{b y d}.
    \end{solu}
\end{exercise}

%========================================================================================
\subsection[Proyecto: la energía en nuestra vida diaria]{Proyecto: formas en las que la energía y sus transformaciones están presentes en nuestra vida diaria}
%========================================================================================

El \textbf{ODS 7} (\emph{Energía asequible y no contaminante}) quiere garantizar a todas las personas el acceso a una energía asequible, segura, sostenible y moderna. Sus metas son el acceso universal a servicios energéticos fiables (7.1), aumentar el peso de las energías renovables (7.2) y duplicar el ritmo de mejora de la eficiencia energética antes de 2030 (7.3).

En este proyecto, que cierra el bloque de la energía, vais a identificar el papel que tienen las distintas formas de energía en vuestra vida diaria y las transformaciones entre ellas, y a mostrarlo en un \textbf{póster} acompañado de un breve reportaje del proceso.

\begin{enumerate}[label=\numberballred]
    \item \textbf{Investigación.} Haced una lista de las formas de energía que han aparecido en las unidades de este bloque: energía cinética, potencial gravitatoria, potencial elástica, mecánica, térmica, química, las ondas sonoras, las ondas electromagnéticas y el Sol como fuente de energía. Formad grupos de tres o cuatro personas y repartid las formas de energía entre los grupos. Buscad información en fuentes fiables sobre cómo está presente vuestra forma de energía en la vida cotidiana, y anotad las fuentes consultadas.
    \item \textbf{Organización.} Diseñad entre toda la clase una ficha común para recoger la información, y seleccionad imágenes propias o con licencia libre (indicando su autoría).
    \item \textbf{Procesamiento y análisis.} Poned en común la información, buscad semejanzas entre las aplicaciones de las distintas formas de energía y estudiad algunas de sus transformaciones (por ejemplo, de potencial a cinética en un salto de agua, de química a térmica en una caldera o de térmica a eléctrica en una central). Debatid cómo se puede contribuir a alcanzar el ODS 7.
    \item \textbf{Elaboración.} Diseñad el póster con la información sobre vuestra forma de energía y sus transformaciones, acompañada de fotografías. Cuidad el diseño: información justa, clara e impactante, y colores y tipografía que se lean bien. Preparad también un reportaje con los momentos más interesantes de la investigación.
    \item \textbf{Comunicación.} Explicad el póster en clase, recoged los comentarios y sugerencias de los compañeros y compañeras para mejorarlo, y exponedlo en el centro, en el aula virtual o en la revista del departamento.
    \item \textbf{Reflexión.} Proponed iniciativas individuales y colectivas: usar energías renovables siempre que sea posible, ahorrar energía, y recordar que, aunque toda la energía mecánica se puede transformar en térmica, la energía térmica no se puede transformar por completo en mecánica (la energía se degrada).
\end{enumerate}
